\fchap{第1章\quad 行列式}

\fsec{深化一　行列式的本质与基本性质}

以下灰底框为高等代数【深化】内容：只补能加深理解、能反哺解题的部分。建议先读正文掌握算法，再回看本段理解“为什么”。

\begin{fdeep}{行列式的本质：反对称多重线性函数}
考研教材里行列式有“六条性质”，一条条背既容易漏，也容易记错方向。
高等代数把这六条压缩成\textbf{三条公理}——凡是满足这三条的映射，有且只有一个，
它就是行列式：
\begin{enumerate}[leftmargin=2.2em,itemsep=0.2em]
\item \textbf{对每一行线性}（多重线性）：固定其余各行时，$|A|$ 是某一行的线性函数；
\item \textbf{换行变号}（反对称性）：交换两行，行列式变号；
\item \textbf{单位阵取 1}（规范性）：$|E|=1$。
\end{enumerate}
\fbref{}：其余性质都能从这三条\textbf{推}出来，不必再背。
\begin{itemize}[leftmargin=2.2em,itemsep=0.2em]
\item 两行相同 $\Rightarrow$ 换这两行后行列式不变又变号 $\Rightarrow 2|A|=0\Rightarrow|A|=0$；
\item 某行全为 $0$ $\Rightarrow$ 由该行线性性得 $|A|=0\cdot(\cdots)=0$；
\item 某行乘 $k$ 加到另一行 $\Rightarrow$ 由线性性拆成“原行列式 $+$ 两行成比例的行列式”$=|A|+0$。
\end{itemize}
\warn{最容易错的一条}：$|kA|=k^{n}|A|$，而\textbf{不是} $k|A|$。
根因就在“对\textbf{每一行}线性”——$n$ 行各出一个 $k$，所以是 $k^{n}$。
把“$k$ 提出来要数行数”记成习惯，这类送分题就不会再错。
\fbref{}服务考纲〔行列式〕“了解行列式的概念，掌握行列式的性质”。把六条性质压缩成\textbf{三条公理}之后，考场上不必回忆性质表；尤其 $|kA|=k^{n}|A|$ 这类“提公因子要数行数”的送分点，根因就在“对每一行线性”。
\end{fdeep}

\begin{fdeep}{几何意义：$|\det A|$ 是体积，符号是定向}
设 $A=(\alpha_1,\dots,\alpha_n)$，把列向量看成 $n$ 个边向量，则
\fml{|\det A|=\text{这 }n\text{ 个列向量张成的平行体的体积}}
而 $\det A$ 的\textbf{符号}记录这组向量相对标准基的\textbf{定向}。
\fml{\det A=0\iff\text{列向量线性相关}\iff\text{平行体“塌缩”到低维}}
\fbref{}：这一条把三个看似无关的章节打通——
行列式为零、向量组线性相关、齐次方程组有非零解、矩阵不可逆、秩 $<n$，
\textbf{说的是同一件事}。做选择题时无论从哪个角度切入都能互相验证，
也解释了为什么“秩”是贯穿全书的中心概念。
\fbref{}服务考纲〔行列式〕与〔矩阵〕。$\det A=0\iff$ 列向量线性相关 $\iff$ 平行体塌缩，把“行列式为零、线性相关、方程组有非零解、不可逆、秩 $<n$”统一成一件事，选择题无论从哪个角度切入都能互相验证。
\end{fdeep}

\begin{fdeep}{行列式项的两种记法与符号的一致性}
容易看出，当行列式的元素全是数域 $P$ 中的数时，它的值也是数域 $P$ 中的一个数.

在行列式的定义中，为了决定每一项的正负号，我们把 $n$ 个元素按行指标排起来.事实上，数的乘法是交换的，因而这 $n$ 个元素的次序是可以任意写的，一般地，$n$ 阶行列式中的项可以写成
\fml{a_{i_1j_1}a_{i_2j_2}\cdots a_{i_nj_n},\qquad (11)}
其中 $i_1i_2\cdots i_n$，$j_1j_2\cdots j_n$ 是两个 $n$ 阶排列.利用排列的性质，不难证明，(11) 的符号等于
\fml{(-1)^{\tau(i_1i_2\cdots i_n)+\tau(j_1j_2\cdots j_n)}.\qquad (12)}

事实上，为了根据定义来决定 (11) 的符号，就要把这 $n$ 个元素重新排一下，使得它们的行指标成自然顺序，也就是排成
\fml{a_{1j_1'}a_{2j_2'}\cdots a_{nj_n'},\qquad (13)}
于是它的符号是
\fml{(-1)^{\tau(j_1'j_2'\cdots j_n')}.\qquad (14)}

现在来证明，(12) 与 (14) 是一致的.我们知道，由 (11) 变到 (13) 可以经过一系列元素的对换来实现.每作一次对换，元素的行指标与列指标所成的排列 $i_1i_2\cdots i_n$ 与 $j_1j_2\cdots j_n$ 就都同时作一次对换，也就是 $\tau(i_1i_2\cdots i_n)$ 与 $\tau(j_1j_2\cdots j_n)$ 同时改变奇偶性，因而它们的和
\fml{\tau(i_1i_2\cdots i_n)+\tau(j_1j_2\cdots j_n)}
的奇偶性不改变.这就是说，对 (11) 作一次元素的对换不改变 (12) 的值.因此，在一系列对换之后有
\fml{(-1)^{\tau(i_1i_2\cdots i_n)+\tau(j_1j_2\cdots j_n)}=(-1)^{\tau(12\cdots n)+\tau(j_1'j_2'\cdots j_n')}=(-1)^{\tau(j_1'j_2'\cdots j_n')}.}
这就证明了 (12) 与 (14) 是一致的.

例如，$a_{21}a_{32}a_{14}a_{43}$ 是 4 阶行列式中一项，$\tau(2314)=2$，$\tau(1243)=1$，于是它的符号应为 $(-1)^{2+1}=-1$.如按行指标排列起来，就是 $a_{14}a_{21}a_{32}a_{43}$，$\tau(4123)=3$，因而它的符号也是 $(-1)^{3}=-1$.
\fbref{}服务考纲〔行列式〕“了解行列式的概念”。考研里用定义算行列式时，最常见的错法是\textbf{只看行指标或只看列指标算逆序数}。本块证明了 $(-1)^{\tau(i_1\cdots i_n)+\tau(j_1\cdots j_n)}=(-1)^{\tau(j_1'\cdots j_n')}$ 二者等价，所以任选一种方式算都对，但不能把两种混着用。
\end{fdeep}

\begin{fdeep}{按列指标排列的定义式与转置行列式}
按 (12) 来决定行列式中每一项的符号的好处在于，行指标与列指标的地位是对称的，因而为了决定每一项的符号，我们同样可以把每一项按列指标排起来，于是定义又可写成
\fml{\begin{vmatrix}
a_{11} & a_{12} & \cdots & a_{1n}\\
a_{21} & a_{22} & \cdots & a_{2n}\\
\vdots & \vdots & & \vdots\\
a_{n1} & a_{n2} & \cdots & a_{nn}
\end{vmatrix}
=\sum_{j_1j_2\cdots j_n}(-1)^{\tau(j_1j_2\cdots j_n)}a_{1j_1}a_{2j_2}\cdots a_{nj_n}.\qquad (15)}

由此即得行列式的下列性质.

\key{性质 1}\quad 行列互换，行列式不变.即
\fml{\begin{vmatrix}
a_{11} & a_{12} & \cdots & a_{1n}\\
a_{21} & a_{22} & \cdots & a_{2n}\\
\vdots & \vdots & & \vdots\\
a_{n1} & a_{n2} & \cdots & a_{nn}
\end{vmatrix}
=\begin{vmatrix}
a_{11} & a_{21} & \cdots & a_{n1}\\
a_{12} & a_{22} & \cdots & a_{n2}\\
\vdots & \vdots & & \vdots\\
a_{1n} & a_{2n} & \cdots & a_{nn}
\end{vmatrix}.\qquad (16)}

\textbf{证明思路}\quad 元素 $a_{ij}$ 在 (16) 的右端位于第 $j$ 行第 $i$ 列，这就是说，$i$ 是它的列指标，$j$ 是它的行指标.因之，把右端按 (15) 展开就等于
\fml{\sum_{j_1j_2\cdots j_n}(-1)^{\tau(j_1j_2\cdots j_n)}a_{1j_1}a_{2j_2}\cdots a_{nj_n},}
它正是左端按 (6) 的展开式.

(16) 中等式右边的行列式称为左边行列式的转置行列式.性质 1 说明行列式转置，值不变.

性质 1 表明，在行列式中行与列的地位是对称的，因之凡是有关行的性质，对列也同样成立.例如，由 (8) 即得下三角形的行列式
\fml{\begin{vmatrix}
a_{11} & 0 & 0 & \cdots & 0\\
a_{21} & a_{22} & 0 & \cdots & 0\\
\vdots & \vdots & \vdots & & \vdots\\
a_{n1} & a_{n2} & a_{n3} & \cdots & a_{nn}
\end{vmatrix}
=a_{11}a_{22}\cdots a_{nn}.\qquad (17)}

下面我们所谈的行列式的性质大多是对行来说的，对于列也有相同的性质，就不重复了.
\fbref{}服务考纲〔行列式〕“掌握行列式的性质”。$|A^{\mathrm T}|=|A|$ 是\textbf{一切性质对行对列都成立}的根据——考纲只写了“按行（列）展开”，但正因这条，凡对行成立的性质对列自动成立，不必把每条性质记两遍。
\end{fdeep}

\begin{fdeep}{性质 4\quad 两行相同，行列式为零}
再根据排列的性质，我们有：

\key{性质 4}\quad 如果行列式中有两行相同，那么行列式为零.所谓两行相同就是说两行的对应元素都相等.

\textbf{证明}\quad 设行列式
\fml{\begin{vmatrix}
a_{11} & a_{12} & \cdots & a_{1n}\\
\vdots & \vdots & & \vdots\\
a_{i1} & a_{i2} & \cdots & a_{in}\\
\vdots & \vdots & & \vdots\\
a_{k1} & a_{k2} & \cdots & a_{kn}\\
\vdots & \vdots & & \vdots\\
a_{n1} & a_{n2} & \cdots & a_{nn}
\end{vmatrix}
=\sum_{j_1\cdots j_n}(-1)^{\tau(j_1\cdots j_{i-1}j_ij_{i+1}\cdots j_{k-1}j_kj_{k+1}\cdots j_n)}
a_{1j_1}\cdots a_{ij_i}\cdots a_{kj_k}\cdots a_{nj_n}\qquad (2)}
中第 $i$ 行与第 $k$ 行相同，即
\fml{a_{ij}=a_{kj},\qquad j=1,2,\cdots,n.\qquad (3)}

为了证明 (2) 为零，只需证明 (2) 的右端所出现的项全能两两相消就行了.事实上，与项
\fml{(-1)^{\tau(j_1\cdots j_i\cdots j_k\cdots j_n)}a_{1j_1}\cdots a_{ij_i}\cdots a_{kj_k}\cdots a_{nj_n}}
同时出现的还有
\fml{(-1)^{\tau(j_1\cdots j_k\cdots j_i\cdots j_n)}a_{1j_1}\cdots a_{ij_k}\cdots a_{kj_i}\cdots a_{nj_n}.}
比较这两项，由 (3) 有
\fml{a_{ij_i}=a_{kj_i},\qquad a_{ij_k}=a_{kj_k}.}
也就是说，这两项有相同的数值.但是排列
\fml{j_1\cdots j_i\cdots j_k\cdots j_n\quad\text{与}\quad j_1\cdots j_k\cdots j_i\cdots j_n}
相差一个对换，因而有相反的奇偶性，所以这两项的符号相反.易知，全部 $n$ 阶排列可以按上述形式两两配对.因之，在 (2) 的右端，对于每一项都有一数值相同但符号相反的项与之成对出现，从而行列式为零.
\fbref{}服务考纲〔行列式〕“掌握行列式的性质”。这条用\textbf{配对相消}证明：把 n! 个排列按“交换第 i,j 个下标”两两配对，每对数值相同、符号相反。理解这个证法后，“两行成比例也为零”可自行推出（先提公因子再归到本性质），不必单记。
\end{fdeep}

\begin{fdeep}{性质 7\quad 对换行列式中两行的位置，行列式反号}
\key{性质 7}\quad 对换行列式中两行的位置，行列式反号.

\textbf{证明}
\fml{\begin{aligned}
&\begin{vmatrix}
a_{11} & a_{12} & \cdots & a_{1n}\\
\vdots & \vdots & & \vdots\\
a_{i1} & a_{i2} & \cdots & a_{in}\\
\vdots & \vdots & & \vdots\\
a_{k1} & a_{k2} & \cdots & a_{kn}\\
\vdots & \vdots & & \vdots\\
a_{n1} & a_{n2} & \cdots & a_{nn}
\end{vmatrix}
=\begin{vmatrix}
a_{11} & a_{12} & \cdots & a_{1n}\\
\vdots & \vdots & & \vdots\\
a_{i1}+a_{k1} & a_{i2}+a_{k2} & \cdots & a_{in}+a_{kn}\\
\vdots & \vdots & & \vdots\\
a_{k1} & a_{k2} & \cdots & a_{kn}\\
\vdots & \vdots & & \vdots\\
a_{n1} & a_{n2} & \cdots & a_{nn}
\end{vmatrix}\\[1mm]
&\quad=\begin{vmatrix}
a_{11} & a_{12} & \cdots & a_{1n}\\
\vdots & \vdots & & \vdots\\
a_{i1}+a_{k1} & a_{i2}+a_{k2} & \cdots & a_{in}+a_{kn}\\
\vdots & \vdots & & \vdots\\
-a_{i1} & -a_{i2} & \cdots & -a_{in}\\
\vdots & \vdots & & \vdots\\
a_{n1} & a_{n2} & \cdots & a_{nn}
\end{vmatrix}\\[1mm]
&\quad=\begin{vmatrix}
a_{11} & a_{12} & \cdots & a_{1n}\\
\vdots & \vdots & & \vdots\\
a_{k1} & a_{k2} & \cdots & a_{kn}\\
\vdots & \vdots & & \vdots\\
-a_{i1} & -a_{i2} & \cdots & -a_{in}\\
\vdots & \vdots & & \vdots\\
a_{n1} & a_{n2} & \cdots & a_{nn}
\end{vmatrix}
=-\begin{vmatrix}
a_{11} & a_{12} & \cdots & a_{1n}\\
\vdots & \vdots & & \vdots\\
a_{k1} & a_{k2} & \cdots & a_{kn}\\
\vdots & \vdots & & \vdots\\
a_{i1} & a_{i2} & \cdots & a_{in}\\
\vdots & \vdots & & \vdots\\
a_{n1} & a_{n2} & \cdots & a_{nn}
\end{vmatrix}.
\end{aligned}}
这里，第一步是把第 $k$ 行加到第 $i$ 行，第二步是把第 $i$ 行的 $-1$ 倍加到第 $k$ 行，第三步是把第 $k$ 行加到第 $i$ 行，最后再把第 $k$ 行的公因子 $-1$ 提出.
\fbref{}服务考纲〔行列式〕“掌握行列式的性质”。这里的三步凑 -1 证明值得记：先加到第 i 行、再把第 i 行的 -1 倍加到第 k 行、最后加回——只用“倍加不变”就推出了“对换反号”。这说明\textbf{倍加变换是最基本的手段}。
\end{fdeep}

\begin{fdeep}{§5 行列式的计算：化上三角形法}
下面我们利用行列式的性质给出一个计算行列式的方法.

在 §3 我们看到，一个上三角形行列式
\fml{\begin{vmatrix}
a_{11} & a_{12} & a_{13} & \cdots & a_{1n}\\
0 & a_{22} & a_{23} & \cdots & a_{2n}\\
0 & 0 & a_{33} & \cdots & a_{3n}\\
\vdots & \vdots & \vdots & & \vdots\\
0 & 0 & 0 & \cdots & a_{nn}
\end{vmatrix}}
就等于它主对角线上元素的乘积
\fml{a_{11}a_{22}\cdots a_{nn}.}
这个计算是很简单的.下面我们想办法把任意的 $n$ 阶行列式化为上三角形行列式来计算.
\fbref{}服务考纲〔行列式〕“会应用行列式的性质和按行（列）展开定理计算行列式”。上三角形行列式等于主对角线乘积，而任意行列式都可用初等行变换化为上三角——这是计算高阶行列式的\textbf{标准算法}，比直接展开快得多。\warn{注意}只用“倍加”与“交换”，因为提公因子会把数字变难算。
\end{fdeep}

\begin{fdeep}{乘法定理 $|AB|=|A||B|$ 与“一证多得”}
这个定理是以下一串结论的\textbf{共同源头}：
\fml{|AB|=|A||B|,\qquad |A^{k}|=|A|^{k},\qquad |A^{-1}|=|A|^{-1}}
\fml{|A^{*}|=|A|^{n-1},\qquad A\sim B\Rightarrow|A|=|B|}
\textbf{证明思路（分块初等变换）}：构造 $2n$ 阶分块矩阵
\fml{M=\begin{pmatrix}E&A\\ B&O\end{pmatrix}}
第一步，作分块初等变换“第二块行 $\leftarrow$ 第二块行 $-B\times$ 第一块行”，
这不改变行列式的值，于是
\fml{M\longrightarrow\begin{pmatrix}E&A\\ O&-BA\end{pmatrix}}
第二步，按 Laplace 定理沿前 $n$ 行展开：
前 $n$ 行中只有左上角的 $E$ 能取到非零的 $n$ 阶子式，
而把 $E$ 所在的行列调到左上角共需 $n^{2}$ 次对换，故有符号因子 $(-1)^{n^{2}}=(-1)^{n}$，于是
\fml{\begin{vmatrix}E&A\\ O&-BA\end{vmatrix}=(-1)^{n}\,|{-BA}|=(-1)^{n}\cdot(-1)^{n}|AB|=|A||B|}
\fbref{}：这里用到的两件事——“分块初等变换不改变行列式”与“Laplace 展开”——
正好是 9 讲里两个看似孤立的技巧，它们在这里合起来一次解决了乘法定理。
\fbref{}：不必记忆五个孤立公式——它们都是同一个定理的推论。
尤其 $A\sim B\Rightarrow|A|=|B|$ 直接来自 $|P^{-1}AP|=|P^{-1}||A||P|=|A|$，
这正是“相似矩阵行列式相等”的根，也是第 8 章判断相似的必要条件之一。
\fbref{}服务考纲〔矩阵〕“了解方阵乘积的行列式的性质”。$|AB|=|A||B|$ 是 $|A^{k}|=|A|^{k}$、$|A^{-1}|=|A|^{-1}$、$|A^{*}|=|A|^{n-1}$、$A\sim B\Rightarrow|A|=|B|$ 的\textbf{共同源头}——不必逐个背，写出本定理再问“我要哪一项”即可。
\end{fdeep}

\fsec{一、具体型行列式的计算：$a_{ij}$ 已给出}

\begin{fcard}{1. 化为基本形行列式}
所谓基本形行列式，是指化至此行列式即可得到结果.

（1）主对角线行列式.
\fmll{\begin{aligned}
&\begin{vmatrix}a_{11}&a_{12}&\cdots&a_{1n}\\0&a_{22}&\cdots&a_{2n}\\\vdots&\vdots&&\vdots\\0&0&\cdots&a_{nn}\end{vmatrix}
=\begin{vmatrix}a_{11}&0&\cdots&0\\a_{21}&a_{22}&\cdots&0\\\vdots&\vdots&&\vdots\\a_{n1}&a_{n2}&\cdots&a_{nn}\end{vmatrix}\\
&=\begin{vmatrix}a_{11}&0&\cdots&0\\0&a_{22}&\cdots&0\\\vdots&\vdots&&\vdots\\0&0&\cdots&a_{nn}\end{vmatrix}=\prod_{i=1}^{n}a_{ii}.
\end{aligned}}

（2）副对角线行列式.
\fmll{\begin{aligned}
&\begin{vmatrix}a_{11}&a_{12}&\cdots&a_{1,n-1}&a_{1n}\\a_{21}&a_{22}&\cdots&a_{2,n-1}&0\\\vdots&\vdots&&\vdots&\vdots\\a_{n1}&0&\cdots&0&0\end{vmatrix}
=\begin{vmatrix}0&\cdots&0&a_{1n}\\0&\cdots&a_{2,n-1}&a_{2n}\\\vdots&&\vdots&\vdots\\a_{n1}&\cdots&a_{n,n-1}&a_{nn}\end{vmatrix}\\
&=\begin{vmatrix}0&\cdots&0&a_{1n}\\0&\cdots&a_{2,n-1}&0\\\vdots&&\vdots&\vdots\\a_{n1}&\cdots&0&0\end{vmatrix}
=(-1)^{\frac{n(n-1)}{2}}a_{1n}a_{2,n-1}\cdots a_{n1}.
\end{aligned}}

（3）拉普拉斯展开式.
设 $A$ 为 $m$ 阶矩阵，$B$ 为 $n$ 阶矩阵，则
\fml{\begin{vmatrix}A&O\\O&B\end{vmatrix}=\begin{vmatrix}A&C\\O&B\end{vmatrix}=\begin{vmatrix}A&O\\C&B\end{vmatrix}=|A||B|,}
\fml{\begin{vmatrix}O&A\\B&O\end{vmatrix}=\begin{vmatrix}C&A\\B&O\end{vmatrix}=\begin{vmatrix}O&A\\B&C\end{vmatrix}=(-1)^{mn}|A||B|.}

（4）范德蒙德行列式.
\fml{\begin{vmatrix}1&1&\cdots&1\\x_1&x_2&\cdots&x_n\\x_1^2&x_2^2&\cdots&x_n^2\\\vdots&\vdots&&\vdots\\x_1^{n-1}&x_2^{n-1}&\cdots&x_n^{n-1}\end{vmatrix}=\prod_{1\le i<j\le n}(x_j-x_i)\,,\ n\ge2.}
\end{fcard}

\begin{fnote}
【注】（1）若所给行列式就是基本形或接近基本形，则直接套公式或经过简单处理化成基本形后套公式.

（2）简单处理的手段.这是具体型行列式解题的要诀

①按零元素多的行或列展开；

②用行列式的性质对差别最小的“对应位置元素”进行处理，尽可能多地化出零元素，再按此行或列展开；

③对于行和或列和相等的情形，将所有列加到第1列或将所有行加到第1行，提出公因式，再用②，等等.

（3）具体型行列式的元素中若含 $x$，则其为 $x$ 的多项式.
\end{fnote}

\begin{fcard}{例1.1}
已知 $f(x)=\begin{vmatrix}2x+1&3&2x+1&1\\2x&-3&4x&-2\\2x+1&2&2x+1&1\\2x&-4&4x&-2\end{vmatrix}$，$g(x)=\begin{vmatrix}2x+1&1&2x+1&3\\5x+1&-2&4x&-3\\0&1&2x+1&2\\2x&-2&4x&-4\end{vmatrix}$，则方程 $f(x)=g(x)$ 的不同的根的个数为\underline{\hspace{2.6em}}.

\end{fcard}

【解】应填2.

\fmll{\begin{aligned}
f(x)&=\begin{vmatrix}2x+1&3&2x+1&1\\2x&-3&4x&-2\\0&-1&0&0\\2x&-4&4x&-2\end{vmatrix}
=\begin{vmatrix}2x+1&3&2x+1&1\\2x&-3&4x&-2\\0&-1&0&0\\0&-1&0&0\end{vmatrix}\\
&=\begin{vmatrix}2x+1&3&2x+1&1\\2x&-3&4x&-2\\0&-1&0&0\\0&0&0&0\end{vmatrix}=0,
\end{aligned}}

\fmll{\begin{aligned}
g(x)&=\begin{vmatrix}2x+1&1&2x+1&3\\9x+3&0&8x+2&3\\-2x-1&0&0&-1\\6x+2&0&8x+2&2\end{vmatrix}
=-\begin{vmatrix}9x+3&8x+2&3\\-2x-1&0&-1\\6x+2&8x+2&2\end{vmatrix}\\
&=-\begin{vmatrix}3x+1&0&1\\-2x-1&0&-1\\6x+2&8x+2&2\end{vmatrix}=-x(8x+2),
\end{aligned}}


所以 $-x(8x+2)=0$，有两个不相等的根.

\begin{fcard}{例1.2}
已知线性方程组 $\begin{cases}ax_1+x_3=1,\\x_1+ax_2+x_3=0,\\x_1+2x_2+ax_3=0,\\ax_1+bx_2=2\end{cases}$ 有解，其中 $a$，$b$ 为常数.若 $\begin{vmatrix}a&0&1\\1&a&1\\1&2&a\end{vmatrix}=4$，则 $\begin{vmatrix}1&a&1\\1&2&a\\a&b&0\end{vmatrix}=$\underline{\hspace{2.6em}}.

\end{fcard}

【分析】$Ax=b$ 有解 $\Leftrightarrow r(A)=r([A,b])=3<4\Rightarrow|[A,b]|=0$.

【解】应填 8.

由于 $\begin{vmatrix}a&0&1\\1&a&1\\1&2&a\end{vmatrix}=4\ne0$，则 $r\left(\begin{pmatrix}a&0&1\\1&a&1\\1&2&a\\a&b&0\end{pmatrix}\right)=3$，又已知方程组有解，则 $r\left(\begin{pmatrix}a&0&1&1\\1&a&1&0\\1&2&a&0\\a&b&0&2\end{pmatrix}\right)=3$，即
\fml{\begin{vmatrix}a&0&1&1\\1&a&1&0\\1&2&a&0\\a&b&0&2\end{vmatrix}=0,}


有
\fml{-\begin{vmatrix}1&a&1\\1&2&a\\a&b&0\end{vmatrix}+2\begin{vmatrix}a&0&1\\1&a&1\\1&2&a\end{vmatrix}=0,}


故
\fml{\begin{vmatrix}1&a&1\\1&2&a\\a&b&0\end{vmatrix}=8.}

\begin{fnote}
【注】本题为什么不用常规思路呢？若本题按“含参数方程组”的常规思路讨论参数 $a$，$b$，最后进行行列式的计算，会十分烦琐，与命题考查的意图南辕北辙.读者可以想到常规的“含参数方程组”求解的是通解或参数，而本题求的是 $\begin{vmatrix}1&a&1\\1&2&a\\a&b&0\end{vmatrix}$ 的值，显然，可以从具体型行列式计算的角度入手.
\end{fnote}

\begin{fcard}{2. 加边法}

对于某些一开始不宜使用“互换”“倍乘”“倍加”性质的行列式，可以考虑使用加边法：$n$ 阶行列式中添加一行、一列升至 $n+1$ 阶行列式.若添加在第1列，且添加的是 $[\,1,\,0,\,\cdots,\,0\,]^{\mathrm{T}}$，则第1行其余元素可以任意添加，行列式的值不变，即
\fml{D_n=\begin{vmatrix}a_{11}&a_{12}&\cdots&a_{1n}\\a_{21}&a_{22}&\cdots&a_{2n}\\\vdots&\vdots&&\vdots\\a_{n1}&a_{n2}&\cdots&a_{nn}\end{vmatrix}=\begin{vmatrix}1&*&*&\cdots&*\\0&a_{11}&a_{12}&\cdots&a_{1n}\\0&a_{21}&a_{22}&\cdots&a_{2n}\\\vdots&\vdots&\vdots&&\vdots\\0&a_{n1}&a_{n2}&\cdots&a_{nn}\end{vmatrix},}

其中 $*$ 处元素可以任意添加.观察原行列式元素的规律性，选择合适的元素填入 $*$ 处，可以使行列式的计算更为简便.
\end{fcard}

\begin{fcard}{例1.3}
设 $\boldsymbol{a}=[\,x_1,\,x_2,\,\cdots,\,x_n\,]^{\mathrm{T}}\ne\boldsymbol{0}$，则 $|\boldsymbol{aa}^{\mathrm{T}}+E|=$\underline{\hspace{2.6em}}.
\end{fcard}

【解】应填 $1+\sum\limits_{i=1}^{n}x_i^2$.

法一
\fmll{\begin{aligned}
|\boldsymbol{aa}^{\mathrm{T}}+E|&=\begin{vmatrix}1+x_1^2&x_1x_2&\cdots&x_1x_n\\x_2x_1&1+x_2^2&\cdots&x_2x_n\\\vdots&\vdots&&\vdots\\x_nx_1&x_nx_2&\cdots&1+x_n^2\end{vmatrix}_n
\overset{(*)}{=}\begin{vmatrix}1&x_1&x_2&\cdots&x_n\\0&1+x_1^2&x_1x_2&\cdots&x_1x_n\\0&x_2x_1&1+x_2^2&\cdots&x_2x_n\\\vdots&\vdots&\vdots&&\vdots\\0&x_nx_1&x_nx_2&\cdots&1+x_n^2\end{vmatrix}_{n+1}\\
&=\begin{vmatrix}1&x_1&x_2&\cdots&x_n\\-x_1&1&0&\cdots&0\\-x_2&0&1&\cdots&0\\\vdots&\vdots&\vdots&&\vdots\\-x_n&0&0&\cdots&1\end{vmatrix}
=\begin{vmatrix}1+\sum\limits_{i=1}^{n}x_i^2&x_1&\cdots&x_n\\0&1&\cdots&0\\\vdots&\vdots&&\vdots\\0&0&\cdots&1\end{vmatrix}
=1+\sum\limits_{i=1}^{n}x_i^2.
\end{aligned}}


法二\ 因为 $\boldsymbol{aa}^{\mathrm{T}}$ 的特征值为 $\sum\limits_{i=1}^{n}x_i^2$，$0$，$0$，$\cdots$，$0$，这里有 $n-1$ 个特征值 $0$，于是 $\boldsymbol{aa}^{\mathrm{T}}+E$ 的特征值为 $1+\sum\limits_{i=1}^{n}x_i^2$，$1$，$1$，$\cdots$，$1$，这里有 $n-1$ 个特征值 $1$.再由第7讲的“一、（2）”知，$|\boldsymbol{aa}^{\mathrm{T}}+E|=\left(1+\sum\limits_{i=1}^{n}x_i^2\right)\times1\times1\times\cdots\times1=1+\sum\limits_{i=1}^{n}x_i^2$.

\begin{fnote}
【注】（$*$）处来自加边法.
\end{fnote}

\begin{fcard}{3. 递推法（高阶$\to$低阶）}

（1）建立递推公式，即建立 $D_n$ 与 $D_{n-1}$ 的关系，有些复杂的题甚至要建立 $D_n$，$D_{n-1}$ 与 $D_{n-2}$ 的关系.

（2）$D_{n-1}$ 与 $D_n$ 要有完全相同的元素分布规律，只是 $D_{n-1}$ 比 $D_n$ 低了一阶.
\end{fcard}

\begin{fcard}{4. 数学归纳法（低阶$\to$高阶）}

涉及 $n$ 阶行列式的\underline{证明型计算问题}，即告知行列式计算结果，让考生证明之，可考虑数学归纳法.

（1）第一数学归纳法（适用于 $F(\,D_n,\,D_{n-1}\,)=0$）：

①验证当 $n=1$ 时，命题成立；

②假设当 $n=k$（$\ge2$）时，命题成立；

③证明当 $n=k+1$ 时，命题成立.

则命题对任意正整数 $n$ 成立.

（2）第二数学归纳法（适用于 $F(\,D_n,\,D_{n-1},\,D_{n-2}\,)=0$）：

①验证当 $n=1$ 和 $n=2$ 时，命题成立；

②假设当 $n<k$ 时，命题成立；

③证明当 $n=k$（$\ge3$）时，命题成立.

则命题对任意正整数 $n$ 成立.
\end{fcard}

\fsec{深化二　余子式、展开定理与行列式计算}

\begin{fdeep}{代数余子式的符号 $(-1)^{i+j}$ 从哪来}
把 $n$ 阶行列式按第 $i$ 行展开：
\fml{|A|=\sum_{j=1}^{n}a_{ij}A_{ij},\qquad A_{ij}=(-1)^{i+j}M_{ij}}
那个 $(-1)^{i+j}$ 不是“规定”，而是\textbf{把 $a_{ij}$ 搬到左上角所付出的换行换列代价}：
要把 $a_{ij}$ 移到 $(1,1)$ 位置，需要把它上方的 $i-1$ 行逐行下移、
左方的 $j-1$ 列逐列右移，共交换 $i+j-2$ 次，故符号为 $(-1)^{i+j-2}=(-1)^{i+j}$。
\fbref{}：符号一旦记错，整题归零，而这是最高频的送分变送命点。
理解来源后在考场上可用“$(1,1)$ 位置是正、棋盘式交错”快速自查。
另外记住\textbf{异行异列配错为零}的恒等式：
\fml{\sum_{k=1}^{n}a_{ik}A_{jk}=\begin{cases}|A|,& i=j\\ 0,& i\ne j\end{cases}}
这正是伴随矩阵 $AA^{*}=|A|E$ 的来源，直接服务第 2 章。
\fbref{}服务考纲〔行列式〕“掌握行列式按行（列）展开定理”。符号一旦记错整题归零，而这是最高频的送分变送命点。理解“换行换列的代价”之后可用“(1,1) 位置为正、棋盘式交错”快速自查；末尾的异行异列配错为零恒等式正是 $AA^{*}=|A|E$ 的来源，直接服务第 2 章。
\end{fdeep}

\begin{fdeep}{余子式与代数余子式（北大 二.6 定义 7、定义 8，印刷页 49、51）}
\key{定义 7}\quad 在行列式
\fml{\begin{vmatrix}
a_{11}&\cdots&a_{1j}&\cdots&a_{1n}\\
\vdots&&\vdots&&\vdots\\
a_{i1}&\cdots&a_{ij}&\cdots&a_{in}\\
\vdots&&\vdots&&\vdots\\
a_{n1}&\cdots&a_{nj}&\cdots&a_{nn}
\end{vmatrix}}
中划去元素 $a_{ij}$ 所在的第 $i$ 行与第 $j$ 列，剩下的 $(n-1)^2$ 个元素按原来的排法构成一个 $n-1$ 阶的行列式
\fml{\begin{vmatrix}
a_{11}&\cdots&a_{1,j-1}&a_{1,j+1}&\cdots&a_{1n}\\
\vdots & & \vdots & \vdots & & \vdots\\
a_{i-1,1}&\cdots&a_{i-1,j-1}&a_{i-1,j+1}&\cdots&a_{i-1,n}\\
a_{i+1,1}&\cdots&a_{i+1,j-1}&a_{i+1,j+1}&\cdots&a_{i+1,n}\\
\vdots & & \vdots & \vdots & & \vdots\\
a_{n1} & \cdots & a_{n,j-1} & a_{n,j+1} & \cdots & a_{nn}
\end{vmatrix}\qquad(3)}
称为元素 $a_{ij}$ 的\textbf{余子式}，记为 $M_{ij}$。

\key{定义 8}\quad 上面所提到的 $A_{ij}$ 称为元素 $a_{ij}$ 的\textbf{代数余子式}，其中
\fml{A_{ij}=(-1)^{i+j}M_{ij}.\qquad(4)}

\fbref{}服务考纲“行列式”：代数余子式是 $n$ 阶行列式降阶与“余子式和代数余子式的计算”的唯一入口，符号 $(-1)^{i+j}$ 必须与位置下标严格对应。
\end{fdeep}

\begin{fdeep}{行列式按一行（列）展开定理（北大 二.6 定理 3，印刷页 49、51）}
对于 $n$ 阶行列式，有
\fml{\begin{vmatrix}
a_{11} & a_{12} & \cdots & a_{1n}\\
\vdots & \vdots & & \vdots\\
a_{i1} & a_{i2} & \cdots & a_{in}\\
\vdots & \vdots & & \vdots\\
a_{n1} & a_{n2} & \cdots & a_{nn}
\end{vmatrix}=a_{i1}A_{i1}+a_{i2}A_{i2}+\cdots+a_{in}A_{in},\quad i=1,2,\cdots,n.\qquad(1)}

三阶行列式可以通过二阶行列式表示：
\fml{\begin{vmatrix}a_{11}&a_{12}&a_{13}\\a_{21}&a_{22}&a_{23}\\a_{31}&a_{32}&a_{33}\end{vmatrix}
=a_{11}\begin{vmatrix}a_{22}&a_{23}\\a_{32}&a_{33}\end{vmatrix}
-a_{12}\begin{vmatrix}a_{21}&a_{23}\\a_{31}&a_{33}\end{vmatrix}
+a_{13}\begin{vmatrix}a_{21}&a_{22}\\a_{31}&a_{32}\end{vmatrix}.\qquad(2)}

这样，公式 $(1)$ 就是说，行列式等于某一行的元素分别与它们代数余子式的乘积之和。在 $(1)$ 中，如果令第 $i$ 行的元素等于另外一行，譬如说，第 $k$ 行的元素，也就是
\fml{a_{ij}=a_{kj},\qquad j=1,\cdots,n,\ k\ne i.}
于是
\fml{a_{k1}A_{i1}+a_{k2}A_{i2}+\cdots+a_{kn}A_{in}
=\begin{vmatrix}
a_{11}&\cdots&a_{1n}\\
\vdots&&\vdots\\
a_{k1}&\cdots&a_{kn}&\text{第 } i \text{ 行}\\
\vdots&&\vdots\\
a_{k1}&\cdots&a_{kn}\\
\vdots&&\vdots\\
a_{n1}&\cdots&a_{nn}
\end{vmatrix}.}
右端的行列式含有两个相同的行，应该为零，这就是说，在行列式中，一行的元素与另一行相应元素的代数余子式的乘积之和为零。

基于行列式中行与列的对称性，在以上的公式和讨论中把行换成列也一样。综上所述，即得

\key{定理 3}\quad 设
\fml{d=\begin{vmatrix}
a_{11} & a_{12} & \cdots & a_{1n}\\
a_{21} & a_{22} & \cdots & a_{2n}\\
\vdots & \vdots & & \vdots\\
a_{n1} & a_{n2} & \cdots & a_{nn}
\end{vmatrix},}
$A_{ij}$ 表示元素 $a_{ij}$ 的代数余子式，则下列公式成立：
\fml{a_{k1}A_{i1}+a_{k2}A_{i2}+\cdots+a_{kn}A_{in}
=\begin{cases}d,&k=i,\\0,&k\ne i;\end{cases}\qquad(6)}
\fml{a_{1l}A_{1j}+a_{2l}A_{2j}+\cdots+a_{nl}A_{nj}
=\begin{cases}d,&l=j,\\0,&l\ne j.\end{cases}\qquad(7)}
用连加号简写为
\fml{\sum_{i=1}^{n}a_{ki}A_{ki}=\begin{cases}d,&k=i,\\0,&k\ne i;\end{cases}\qquad
\sum_{j=1}^{n}a_{lj}A_{lj}=\begin{cases}d,&l=j,\\0,&l\ne j.\end{cases}}

在计算数字行列式时，直接应用展开式 $(6)$ 或 $(7)$ 不一定能简化计算，因为把一个 $n$ 阶行列式的计算换成 $n$ 个 $n-1$ 阶行列式的计算并不减少计算量，只是在行列式中某一行或某一列含较多的零时，应用公式 $(6)$ 或 $(7)$ 才有意义。但这两个公式在理论上是很重要的。

\fbref{}服务考纲“行列式”：这是“降阶法”的定理依据，也是“某行元素乘另一行代数余子式之和为 0”这一常考结论的来源。
\end{fdeep}

\begin{fdeep}{$n=3$ 时展开公式的几何意义（北大 二.6，印刷页 52）}
在 $n=3$ 时，公式 $(6)$ 有明显的几何意义。如果把行列式的行看作向量在直角坐标系下的坐标，即设
\fml{\alpha_1=(a_{11},a_{12},a_{13}),\qquad \alpha_2=(a_{21},a_{22},a_{23}),\qquad \alpha_3=(a_{31},a_{32},a_{33}),}
那么
\fml{\alpha_2\times\alpha_3=(A_{11},A_{12},A_{13}).}
于是
\fml{a_{11}A_{11}+a_{12}A_{12}+a_{13}A_{13}=\alpha_1\cdot(\alpha_2\times\alpha_3),}
\fml{a_{21}A_{11}+a_{22}A_{12}+a_{23}A_{13}=\alpha_2\cdot(\alpha_2\times\alpha_3)=0,}
\fml{a_{31}A_{11}+a_{32}A_{12}+a_{33}A_{13}=\alpha_3\cdot(\alpha_2\times\alpha_3)=0.}

\fbref{}服务考纲“行列式”：给“异行（列）代数余子式之和为 0”一个几何图像——混合积中有两个向量相同即为 0。
\end{fdeep}

\begin{fdeep}{Vandermonde 行列式与缺项型变式}
$n$ 阶 Vandermonde 行列式
\fml{V_{n}=\begin{vmatrix}1&1&\cdots&1\\ x_{1}&x_{2}&\cdots&x_{n}\\ \vdots&\vdots&&\vdots\\ x_{1}^{n-1}&x_{2}^{n-1}&\cdots&x_{n}^{n-1}\end{vmatrix}=\prod_{1\le i<j\le n}(x_{j}-x_{i})}
\textbf{缺项型}：若第一行不是全 $1$，而是 $x_{1}^{k},x_{2}^{k},\dots$（缺了某一次幂），
则可利用\textbf{把缺的那一行补上并作差}，或直接用“按列线性性拆项 + 行列式为零项剔掉”的办法处理，
例如含 $x_{i}^{n}$ 与 $x_{i}^{n-2}$ 而缺 $x_{i}^{n-1}$ 的情形，可用
\fml{\begin{vmatrix}1&1&1\\ x_{1}&x_{2}&x_{3}\\ x_{1}^{3}&x_{2}^{3}&x_{3}^{3}\end{vmatrix}=(x_{1}+x_{2}+x_{3})\prod_{i<j}(x_{j}-x_{i})}
这类结论\textbf{不要死记}，掌握“补行相减 / 拆项消零”两条通道即可现场推出。
\fbref{}：考研高频考 Vandermonde 及其变式，考场上现推比背公式更可靠，
且能避免把缺项型的系数记错。
\end{fdeep}

\begin{fdeep}{各类降阶技巧为什么都有效}
9 讲给出了加边法、递推法、拆项法、升阶法等套路。它们并非各自独立的技巧，
而是\textbf{多重线性}这同一原理的不同用法：
\begin{itemize}[leftmargin=2.2em,itemsep=0.25em]
\item \textbf{拆项法}——直接使用“对某一行线性”：把一行拆成两行之和，行列式拆成两个行列式之和；
\item \textbf{递推法}——利用行列式结构沿阶数下降，同一线性关系在相邻阶之间重复出现；
\item \textbf{加边（升阶）法}——把 $n$ 阶“升”成 $n+1$ 阶，使新行列式出现可利用的整行/整列公因子，
      再用一次行列式性质把多出来的边“消”掉，本质是\textbf{先用规范性固定尺度}；
\item \textbf{“打洞”消元}——用“某行乘 $k$ 加到另一行不改变行列式”，制造大量零元以启用展开。
\end{itemize}
\fbref{}：遇到没见过的高阶行列式时，不要盲目试套路，先问自己
“这个行列式对哪一行、哪一列具有可利用的线性结构”，
四个方法中总有一个能对上号。这比背题型更能应付新题。
\end{fdeep}

\begin{fcard}{例1.4}
$D_n=\begin{vmatrix}b&-1&0&\cdots&0&0\\0&b&-1&\cdots&0&0\\\vdots&\vdots&\vdots&&\vdots&\vdots\\0&0&0&\cdots&b&-1\\a_n&a_{n-1}&a_{n-2}&\cdots&a_2&b+a_1\end{vmatrix}=$\underline{\hspace{2.6em}}.

（异爪形行列式）
\end{fcard}

【解】应填 $b^n+a_1b^{n-1}+a_2b^{n-2}+\cdots+a_{n-1}b+a_n$.

递推法.按第1列展开，得
\fmll{\begin{aligned}
D_n&=b\begin{vmatrix}b&-1&0&\cdots&0&0\\0&b&-1&\cdots&0&0\\\vdots&\vdots&\vdots&&\vdots&\vdots\\a_{n-1}&a_{n-2}&a_{n-3}&\cdots&a_2&b+a_1\end{vmatrix}_{n-1}
+(-1)^{n+1}a_n\begin{vmatrix}-1&0&\cdots&0&0\\b&-1&\cdots&0&0\\\vdots&\vdots&&\vdots&\vdots\\0&0&\cdots&b&-1\end{vmatrix}_{n-1}\\
&=bD_{n-1}+a_n.
\end{aligned}}


下面做递推，得
\fmll{\begin{aligned}
D_n&=bD_{n-1}+a_n=b\left(bD_{n-2}+a_{n-1}\right)+a_n=b^2D_{n-2}+a_{n-1}b+a_n\\
&=b^2\left(bD_{n-3}+a_{n-2}\right)+a_{n-1}b+a_n\\
&=\cdots=b^{n-1}D_1+a_2b^{n-2}+\cdots+a_{n-1}b+a_n,
\end{aligned}}

其中 $D_1\overset{(*)}{=}b+a_1$，故 $D_n=b^n+a_1b^{n-1}+a_2b^{n-2}+\cdots+a_{n-1}b+a_n$.

\begin{fnote}
【注】（1）（$*$）处提醒考生注意，$D_n$ 的元素分布规律应从右下角往左上看，写出 $D_k$（$k=1,\,2,\,\cdots,\,n-1$，$n$）供考生参考：
\fml{D_k=\begin{vmatrix}b&-1&0&\cdots&0&0\\0&b&-1&\cdots&0&0\\\vdots&\vdots&\vdots&&\vdots&\vdots\\0&0&0&\cdots&b&-1\\a_k&a_{k-1}&a_{k-2}&\cdots&a_2&b+a_1\end{vmatrix}_k}

事实上，选第1列展开是基于 $D_n$ 的这种元素分布规律，若选第 $n$ 列展开，余子式便不是 $D_{n-1}$，破坏了元素分布规律，无法建立递推公式.

（2）本题也可用第一数学归纳法做.由
\fml{D_1=b+a_1,}
\fml{D_2=\begin{vmatrix}b&-1\\a_2&b+a_1\end{vmatrix}=b^2+a_1b+a_2,}
\end{fnote}


\begin{fcard}{例 1.5}
证明：$n$ 阶行列式
\fml{D_n=\begin{vmatrix}
2a & 1 & 0 & \cdots & 0 & 0\\
a^{2} & 2a & 1 & \cdots & 0 & 0\\
0 & a^{2} & 2a & \cdots & 0 & 0\\
\vdots & \vdots & \vdots & & \vdots & \vdots\\
0 & 0 & 0 & \cdots & 2a & 1\\
0 & 0 & 0 & \cdots & a^{2} & 2a
\end{vmatrix}=(n+1)a^{n}.}

【证】用第二数学归纳法.

当 $n=1$ 时，$D_1=2a=(1+1)a^{1}$，命题成立.

当 $n=2$ 时，$D_2=\begin{vmatrix}2a & 1\\ a^{2} & 2a\end{vmatrix}=4a^{2}-a^{2}=3a^{2}=(2+1)a^{2}$，命题成立.

假设当 $n<k$ 时，命题成立，则当 $n=k$（$\geqslant 3$）时，$D_k$ 按第 1 列展开，得
\fml{D_k=2aD_{k-1}+(-1)^{1+2}a^{2}\begin{vmatrix}
1 & 0 & 0 & \cdots & 0 & 0\\
a^{2} & 2a & 1 & \cdots & 0 & 0\\
0 & a^{2} & 2a & \cdots & 0 & 0\\
\vdots & \vdots & \vdots & & \vdots & \vdots\\
0 & 0 & 0 & \cdots & 2a & 1\\
0 & 0 & 0 & \cdots & a^{2} & 2a
\end{vmatrix}_{k-1}}
\fml{=2aD_{k-1}-a^{2}D_{k-2}}
\fml{=2a\,(k+1-1)\,a^{k-1}-a^{2}(k-2+1)\,a^{k-2}=(k+1)a^{k},}
得证，命题成立.
\end{fcard}

\begin{fnote}
【注】一般来说，当命题直接要求计算行列式时，优先考虑递推法，如例 1.4；当命题给出行列式的结果，要求证明时，优先考虑数学归纳法，如例 1.5.
\end{fnote}

% -----------------------------------------------------------
% PDF p18（印刷页 7）
% -----------------------------------------------------------
\fsec{二、抽象型行列式的计算：$a_{ij}$ 未给出}

\begin{fcard}{1. 用行列式的性质}
用行列式的性质将所求行列式进一步化成已知行列式.
\end{fcard}

\begin{fcard}{2. 用矩阵知识}
（1）设 $C=AB$，$A$，$B$ 为同阶方阵，则 $|C|=|AB|=|A||B|$.
\end{fcard}

% OPD 蓝色标注（原稿）：→$D_{23}$（化归经典形式）
\begin{fnote}
【注】要善于发现 $AB$，如 $A=[\alpha_1,\alpha_2,\alpha_3]$，$B=\begin{bmatrix}a & b & b\\ b & a & b\\ b & b & a\end{bmatrix}$，$C=\begin{bmatrix}1 & 2 & 3\\ 2 & 3 & 4\\ 3 & 4 & 5\end{bmatrix}$.

令
\fml{D_1=[a\alpha_1+b(\alpha_2+\alpha_3),\,a\alpha_2+b(\alpha_1+\alpha_3),\,a\alpha_3+b(\alpha_1+\alpha_2)]=AB;}
\fml{D_2=[\alpha_1+2\alpha_2+3\alpha_3,\,2\alpha_1+3\alpha_2+4\alpha_3,\,3\alpha_1+4\alpha_2+5\alpha_3]=AC.}
这样利于使用 $|D_1|=|A||B|$，$|D_2|=|A||C|$.
\end{fnote}

（2）设 $C=A+B$，$A$，$B$ 为同阶方阵，则 $|C|=|A+B|$，但由于 $|A+B|$ 不一定等于 $|A|+|B|$，故需对 $|A+B|$ 作恒等变形，转化为矩阵乘积的行列式. 这里的恒等变形一般是：①由题设条件，如 $E=AA^{T}$；②用 $E=AA^{-1}$ 等.
% OPD 蓝色标注（原稿）：→$D_{23}$（化归经典形式）  （指向「转化为矩阵乘积的行列式」）
% OPD 蓝色标注（原稿）：→$D_{22}$（转换等价表述）  （指向「①由题设条件，如 $E=AA^{T}$」）



\fsec{深化三　克拉默法则与方程组的行列式判据}

\begin{fdeep}{分块三角行列式 $\begin{vmatrix}A&0\\C&B\end{vmatrix}=|A|\,|B|$（北大 二.6 例 3，印刷页 54、55）}
\fml{\begin{vmatrix}
a_{11}&\cdots&a_{1k}&0&\cdots&0\\
\vdots & & \vdots & \vdots & & \vdots\\
a_{k1}&\cdots&a_{kk}&0&\cdots&0\\
c_{11}&\cdots&c_{1k}&b_{11}&\cdots&b_{1r}\\
\vdots & & \vdots & \vdots & & \vdots\\
c_{r1}&\cdots&c_{rk}&b_{r1}&\cdots&b_{rr}
\end{vmatrix}
=\begin{vmatrix}
a_{11}&\cdots&a_{1k}\\
\vdots&&\vdots\\
a_{k1}&\cdots&a_{kk}
\end{vmatrix}
\begin{vmatrix}
b_{11}&\cdots&b_{1r}\\
\vdots&&\vdots\\
b_{r1}&\cdots&b_{rr}
\end{vmatrix}.\qquad(9)}
证明思路：对 $k$ 用数学归纳法。当 $k=1$ 时，$(9)$ 的左端为
\fml{\begin{vmatrix}
a_{11}&0&\cdots&0\\
c_{11}&b_{11}&\cdots&b_{1r}\\
\vdots & \vdots & & \vdots\\
c_{r1}&b_{r1}&\cdots&b_{rr}
\end{vmatrix},}
按第一行展开，就得到所要的结论。

假设 $(9)$ 对 $k=m-1$ 时（即左端行列式的左上角是 $m-1$ 阶时）已经成立，现在来看 $k=m$ 的情形，按第一行展开，有
\fml{\begin{vmatrix}
a_{11}&\cdots&a_{1m}&0&\cdots&0\\
\vdots & & \vdots & \vdots & & \vdots\\
a_{m1}&\cdots&a_{mm}&0&\cdots&0\\
c_{11}&\cdots&c_{1m}&b_{11}&\cdots&b_{1r}\\
\vdots & & \vdots & \vdots & & \vdots\\
c_{r1}&\cdots&c_{rm}&b_{r1}&\cdots&b_{rr}
\end{vmatrix}
=a_{11}\begin{vmatrix}
a_{22}&\cdots&a_{2m}&0&\cdots&0\\
\vdots & & \vdots & \vdots & & \vdots\\
a_{m2}&\cdots&a_{mm}&0&\cdots&0\\
c_{12}&\cdots&c_{1m}&b_{11}&\cdots&b_{1r}\\
\vdots & & \vdots & \vdots & & \vdots\\
c_{r2}&\cdots&c_{rm}&b_{r1}&\cdots&b_{rr}
\end{vmatrix}+\cdots}
\fml{\cdots+(-1)^{1+m}a_{1m}
\begin{vmatrix}
a_{21}&\cdots&a_{2,m-1}\\
\vdots&&\vdots\\
a_{m1}&\cdots&a_{m,m-1}
\end{vmatrix}
\begin{vmatrix}
b_{11}&\cdots&b_{1r}\\
\vdots&&\vdots\\
b_{r1}&\cdots&b_{rr}
\end{vmatrix}
=\begin{vmatrix}
a_{11}&\cdots&a_{1m}\\
\vdots&&\vdots\\
a_{m1}&\cdots&a_{mm}
\end{vmatrix}
\begin{vmatrix}
b_{11}&\cdots&b_{1r}\\
\vdots&&\vdots\\
b_{r1}&\cdots&b_{rr}
\end{vmatrix}.}
这里第二个等号是用了归纳法假定，最后一步是根据按一行展开的公式。根据归纳法原理，$(9)$ 式普遍成立。

\fbref{}服务考纲“行列式”“矩阵”：分块上（下）三角行列式 $=|A||B|$ 是含零块矩阵行列式、$\begin{vmatrix}A&O\\O&B\end{vmatrix}=|A||B|$ 类考题的直接依据，也是拉普拉斯定理最常见的特例。
\end{fdeep}

\begin{fdeep}{克拉默（Cramer）法则：定理 4 的陈述（北大 二.7，印刷页 55）}
现在我们来应用行列式解决线性方程组的问题。在这里只考虑方程个数与未知量的个数相等的情形。以后会看到，这是一个重要的情形。至于更一般的情形留到下一章讨论。

\key{定理 4}\quad 如果线性方程组
\fml{\begin{cases}
a_{11}x_1+a_{12}x_2+\cdots+a_{1n}x_n=b_1,\\
a_{21}x_1+a_{22}x_2+\cdots+a_{2n}x_n=b_2,\\
\qquad\cdots\cdots\cdots\cdots\cdots\cdots\\
a_{n1}x_1+a_{n2}x_2+\cdots+a_{nn}x_n=b_n
\end{cases}\qquad(1)}
的系数矩阵
\fml{A=\begin{pmatrix}
a_{11} & a_{12} & \cdots & a_{1n}\\
a_{21} & a_{22} & \cdots & a_{2n}\\
\vdots & \vdots & & \vdots\\
a_{n1} & a_{n2} & \cdots & a_{nn}
\end{pmatrix}\qquad(2)}
的行列式，即系数行列式
\fml{d=|A|\ne 0,}
那么线性方程组 $(1)$ 有解，并且解是唯一的，解可以通过系数表为
\fml{x_1=\frac{d_1}{d},\quad x_2=\frac{d_2}{d},\quad \cdots,\quad x_n=\frac{d_n}{d},\qquad(3)}
其中 $d_j$ 是把矩阵 $A$ 中第 $j$ 列换成方程组的常数项 $b_1,b_2,\cdots,b_n$ 所成的矩阵的行列式，即
\fml{d_j=\begin{vmatrix}
a_{11}&\cdots&a_{1,j-1}&b_1&a_{1,j+1}&\cdots&a_{1n}\\
a_{21}&\cdots&a_{2,j-1}&b_2&a_{2,j+1}&\cdots&a_{2n}\\
\vdots & & \vdots & \vdots & \vdots & & \vdots\\
a_{n1}&\cdots&a_{n,j-1}&b_n&a_{n,j+1}&\cdots&a_{nn}
\end{vmatrix},\qquad j=1,2,\cdots,n.\qquad(4)}

定理中包含着三个结论：$1^{\circ}$ 方程组有解；$2^{\circ}$ 解是唯一的；$3^{\circ}$ 解由公式 $(3)$ 给出。这三个结论是有联系的，因此证明的步骤是：
1.\ 把 $\left(\dfrac{d_1}{d},\dfrac{d_2}{d},\cdots,\dfrac{d_n}{d}\right)$ 代入方程组，验证它确是解。

2.\ 假如方程组有解，证明它的解必由公式 $(3)$ 给出。

\fbref{}服务考纲“线性方程组”：Cramer 法则给出“解 $=$ 行列式之比”的显式公式，是判断解的唯一性、构造反例（$|A|=0$ 时失效）与含参数方程组讨论的出发点。
\end{fdeep}

\begin{fdeep}{克拉默法则的证明思路：用代数余子式导出解（北大 二.7，印刷页 56、57）}
在下面的证明中，为了写起来简短些，我们将尽量用连加号 $\sum$。

\textbf{证明}\quad 1.\ 把方程组 $(1)$ 简写为
\fml{\sum_{j=1}^{n}a_{ij}x_j=b_i,\qquad i=1,2,\cdots,n.\qquad(5)}
首先来证明 $(3)$ 的确是 $(1)$ 的解。把 $(3)$ 代入第 $i$ 个方程，左端为
\fml{\sum_{j=1}^{n}a_{ij}\frac{d_j}{d}=\frac{1}{d}\sum_{j=1}^{n}a_{ij}d_j.\qquad(6)}
因为
\fml{d_j=b_1A_{1j}+b_2A_{2j}+\cdots+b_nA_{nj}=\sum_{s=1}^{n}b_sA_{sj},}
所以
\fml{\frac{1}{d}\sum_{j=1}^{n}a_{ij}d_j=\frac{1}{d}\sum_{j=1}^{n}a_{ij}\sum_{s=1}^{n}b_sA_{sj}
=\frac{1}{d}\sum_{s=1}^{n}\sum_{j=1}^{n}a_{ij}A_{sj}b_s
=\frac{1}{d}\sum_{s=1}^{n}\left(\sum_{j=1}^{n}a_{ij}A_{sj}\right)b_s.}
根据定理 $3$ 中公式 $(6)$，有
\fml{\frac{1}{d}\left(\sum_{j=1}^{n}a_{ij}A_{ij}\right)b_i=\frac{1}{d}\cdot db_i=b_i.}
这与第 $i$ 个方程的右端一致。换句话说，把 $(3)$ 代入方程使它们同时变成恒等式，因而 $(3)$ 确为方程组 $(1)$ 的解。

2.\ 设 $(c_1,c_2,\cdots,c_n)$ 是方程组 $(1)$ 的一个解，于是有 $n$ 个恒等式
\fml{\sum_{j=1}^{n}a_{ij}c_j=b_i,\qquad i=1,2,\cdots,n.\qquad(7)}
为了证明 $c_k=\dfrac{d_k}{d}$，我们取系数矩阵中第 $k$ 列元素的代数余子式 $A_{1k},A_{2k},\cdots,A_{nk}$，用它们分别乘 $(7)$ 中 $n$ 个恒等式，有
\fml{A_{ik}\sum_{j=1}^{n}a_{ij}c_j=b_iA_{ik},\qquad i=1,2,\cdots,n,}
这仍是 $n$ 个恒等式。把它们加起来，即得
\fml{\sum_{i=1}^{n}A_{ik}\sum_{j=1}^{n}a_{ij}c_j=\sum_{i=1}^{n}b_iA_{ik}.\qquad(8)}
等式右端等于在行列式 $d$ 按第 $k$ 列的展开式中把 $a_{ik}$ 分别换成 $b_i\ (i=1,2,\cdots,n)$，因此，它等于把行列式 $d$ 中第 $k$ 列换成 $b_1,b_2,\cdots,b_n$ 所得的行列式，也就是 $d_k$。再来看 $(8)$ 的左端，即
\fml{\sum_{i=1}^{n}A_{ik}\sum_{j=1}^{n}a_{ij}c_j=\sum_{i=1}^{n}\sum_{j=1}^{n}a_{ij}A_{ik}c_j
=\sum_{j=1}^{n}\left(\sum_{i=1}^{n}a_{ij}A_{ik}\right)c_j.}
由上节定理 $3$ 中公式 $(7)$，
\fml{\sum_{i=1}^{n}a_{ij}A_{ik}=\begin{cases}d,&j=k,\\0,&j\ne k.\end{cases}}
所以
\fml{\sum_{j=1}^{n}\left(\sum_{i=1}^{n}a_{ij}A_{ik}\right)c_j=dc_k.}
于是，$(8)$ 即为
\fml{dc_k=d_k,\qquad k=1,2,\cdots,n,}
也就是说
\fml{c_k=\frac{d_k}{d},\qquad k=1,2,\cdots,n.}
这就是说，如果 $(c_1,c_2,\cdots,c_n)$ 是方程组的一个解，它必为
\fml{\left(\frac{d_1}{d},\frac{d_2}{d},\cdots,\frac{d_n}{d}\right),\qquad(9)}
因而方程组最多有一组解。

定理 $4$ 通常称为\textbf{克拉默法则}。

\fbref{}服务考纲“线性方程组”：证明的两步恰好是两条常考思路——“用代数余子式乘方程再相加”这一技巧，以及“把 $d_j$ 看成第 $j$ 列被替换后的展开式”这一恒等变形。
\end{fdeep}

\begin{fdeep}{克拉默法则的适用边界与意义（北大 二.7，印刷页 58）}
应该注意，定理 $4$ 所讨论的只是系数矩阵的行列式不为零的方程组，它只能应用于这种方程组；至于方程组的系数行列式为零的情形，将在下一章的一般情形中一并讨论。

克拉默法则的意义主要在于它给出了解与系数的明显关系，这一点在以后许多问题的讨论中是重要的。但是用克拉默法则进行计算是不方便的，因为按这一法则解一个 $n$ 个未知量 $n$ 个方程的线性方程组就要计算 $n+1$ 个 $n$ 阶行列式，这个计算量是很大的。

\fbref{}服务考纲“线性方程组”：明确“方阵且 $|A|\ne 0$”这一前置条件，避免对方程数 $\ne$ 未知数或 $|A|=0$ 的方程组误用 Cramer 法则；并说明它是理论工具而非计算工具。
\end{fdeep}

\begin{fdeep}{齐次线性方程组与 $|A|=0$（北大 二.7 定理 5，印刷页 58）}
常数项全为零的线性方程组称为\textbf{齐次线性方程组}。显然，齐次线性方程组总是有解的，因为 $(0,0,\cdots,0)$ 就是一个解，它称为\textbf{零解}。对于齐次线性方程组，我们关心的问题常常是，它除去零解以外还有没有其他解，或者说，它有没有\textbf{非零解}。对于方程个数与未知量个数相同的齐次线性方程组，应用克拉默法则就有

\key{定理 5}\quad 如果齐次线性方程组
\fml{\begin{cases}
a_{11}x_1+a_{12}x_2+\cdots+a_{1n}x_n=0,\\
a_{21}x_1+a_{22}x_2+\cdots+a_{2n}x_n=0,\\
\qquad\cdots\cdots\cdots\cdots\cdots\cdots\\
a_{n1}x_1+a_{n2}x_2+\cdots+a_{nn}x_n=0
\end{cases}\qquad(10)}
的系数矩阵的行列式 $|A|\ne 0$，那么它只有零解。换句话说，如果方程组 $(10)$ 有非零解，那么必有 $|A|=0$。

\textbf{证明}\quad 应用克拉默法则，因为行列式 $d_j$ 中有一列为零，所以
\fml{d_j=0,\qquad j=1,2,\cdots,n.}
这就是说，它的唯一的解是
\fml{\left(\frac{d_1}{d},\frac{d_2}{d},\cdots,\frac{d_n}{d}\right)=(0,0,\cdots,0).}

\fbref{}服务考纲“线性方程组”：$n$ 个方程 $n$ 个未知量的齐次方程组“有非零解 $\Leftrightarrow |A|=0$”是特征值（$|A-\lambda E|=0$）与线性相关判定的公共接口。
\end{fdeep}

\begin{fcard}{例 1.6}
设 $A$，$B$ 为 $n$ 阶正交矩阵，则 $n$ 为奇数时，$|(A-B)(A+B)|=\underline{\qquad}$.

【分析】$AA^{T}=E(A^{T}=A^{-1})$；$BB^{T}=E(B^{T}=B^{-1})$.

【解】应填 0.
% OPD 蓝色标注（原稿）：=$E$（$D_{22}$（转换等价表述））
\fml{|(A-B)(A+B)|=|(A-B)^{T}(A+B)|=|A^{T}A-B^{T}A+A^{T}B-B^{T}B|=|A^{T}B-B^{T}A|,}
又 $|(A-B)(A+B)|=|(A+B)^{T}(A-B)|=(-1)^{n}|A^{T}B-B^{T}A|$，故当 $n$ 为奇数时，$|(A-B)(A+B)|=0$.
\end{fcard}

\begin{fnote}
【注】$|A|=0$ 可由以下条件推知.

①$|A|=k|A|$，$k\neq 1$，常考 $|A|=(-1)^{n}|A|$($n$ 为奇数).

②$A$ 有特征值 0.

③$A_{n\times n}x=\mathbf{0}$ 有非零解.

④$r(A_{m\times n})<n$.

⑤$A$ 的极大无关组成员个数 $<n$（也即 $A$ 的行（列）向量组线性相关）.

⑥实在不行，或显而易见时，立即推——放弃考研！不，用反证法，设 $A$ 可逆，即 $|A|\neq 0$，推得矛盾，即可得证！
\end{fnote}

% -----------------------------------------------------------
% PDF p19（印刷页 8）
% -----------------------------------------------------------
% OPD 蓝色标注（原稿）：→$D_1$（常规操作）$+D_{23}$（化归经典形式）
（3）设 $A$ 为 $n$ 阶矩阵，则 $|A^{*}|=|A|^{n-1}$，$|(A^{*})^{*}|=||A|^{n-2}A|=|A|^{(n-1)^{2}}$. 更全面的公式总结在第 3 讲的“四、2.（2）”处.

\begin{fcard}{例 1.7}
设 $A$ 是 4 阶矩阵，且其伴随矩阵 $A^{*}$ 的特征值为 1，$-1$，$-2$，4，则 $|A^{3}+2A^{2}-A-3E|=\underline{\qquad}$.

【解】应填 $-\dfrac{253}{8}$.

% 原稿此处右侧另有一张 3 行 3 列表格
\fml{\begin{array}{|c|c|c|}\hline
A & A^{*} & f(A)\\ \hline
\lambda & \dfrac{|A|}{\lambda} & f(\lambda)\\ \hline
\xi & \xi & \xi\\ \hline
\end{array}}

由于 $|A^{*}|=1\times(-1)\times(-2)\times4=8\neq 0$，可知 $A^{*}$ 可逆，于是 $A$ 可逆. 又 $|A^{*}|=|A|^{n-1}=|A|^{3}=8$，得 $|A|=2$. 故 $A$ 的特征值 $\lambda_A=\dfrac{|A|}{\lambda_{A^{*}}}$，即为 $2$，$-2$，$-1$，$\dfrac{1}{2}$.

设 $f(A)=A^{3}+2A^{2}-A-3E$，则 $f(A)$ 的特征值为

% 原稿蓝色旁注（交叉引用）：见第 7 讲的“一、（3）①”中的表格：$f(A)$ 的特征值为$f(\lambda)$
见第 7 讲的“一、（3）①”中的表格：$f(A)$ 的特征值为 $f(\lambda)$.

\fml{f(2)=2^{3}+2\times2^{2}-2-3=11,}
\fml{f(-2)=(-2)^{3}+2\times(-2)^{2}+2-3=-1,}
\fml{f(-1)=-1+2+1-3=-1,}
\fml{f\!\left(\frac{1}{2}\right)=\left(\frac{1}{2}\right)^{3}+2\times\left(\frac{1}{2}\right)^{2}-\frac{1}{2}-3=-\frac{23}{8},}
故
\fml{|A^{3}+2A^{2}-A-3E|=f(2)\cdot f(-2)\cdot f(-1)\cdot f\!\left(\frac{1}{2}\right)=-\frac{253}{8}.}
\end{fcard}

\begin{fcard}{3. 用相似理论}
% OPD 蓝色标注（原稿）：→$D_1$（常规操作）$+D_{22}$（转换等价表述）
若 $A$ 相似于 $B$，则 $|A|=|B|$，且 $|A|=\prod\limits_{i=1}^{n}\lambda_i$，主要考：

（1）用行列式 $|\lambda E-A|=0$ 求特征值；

（2）用特征值求行列式.
\end{fcard}

\begin{fcard}{例 1.8}
设 $A=k\begin{bmatrix}
1 & a & a & \cdots & a\\
a & 1 & a & \cdots & a\\
\vdots & \vdots & \vdots & & \vdots\\
a & a & a & \cdots & 1
\end{bmatrix}_{n\times n}$，$0<a\leqslant 1$，$k>0$，求 $A$ 的最大特征值.

% ---------------------------------------------------------
% PDF p20（印刷页 9）
% ---------------------------------------------------------
【解】
\fmll{|\lambda E-A|=\begin{vmatrix}
\lambda-k & -ka & \cdots & -ka\\
-ka & \lambda-k & \cdots & -ka\\
\vdots & \vdots & & \vdots\\
-ka & -ka & \cdots & \lambda-k
\end{vmatrix}=[\lambda-k+(n-1)(-ka)](\lambda-k+ka)^{n-1}=0,}

% 原稿蓝色旁注（带箭头指向上式结果）：用好公式：
用好公式：$\begin{bmatrix}a & b & \cdots & b\\ b & a & \cdots & b\\ \vdots & \vdots & & \vdots\\ b & b & \cdots & a\end{bmatrix}=[a+(n-1)b](a-b)^{n-1}$

故 $\lambda_1=k+(n-1)ka=k[1+(n-1)a]$（单重），$\lambda_2=k-ka=k(1-a)$（$n-1$ 重）.

由 $1+(n-1)a>1-a$，$k>0$，故最大特征值为 $k[1+(n-1)a]$.
\end{fcard}
